% arXiv submission — A perturbation ladder for SAE dictionaries
% J. Melton, American Code Labs — September 2026
%
% DRAFT. Unmeasured values are \pending{} and render red; the banner counts them.
%
% Framing note: this paper is a direct complement to Duan (arXiv:2606.03002), which
% asks whether a FROZEN dictionary's features still fire the same way under
% quantization. We ask whether the dictionary you would LEARN changes. The two
% metrics are not comparable numerically and the draft must never imply they are.

\documentclass[11pt,a4paper]{article}

\usepackage[utf8]{inputenc}
\usepackage[T1]{fontenc}
\usepackage{lmodern}
\usepackage{amsmath}
\usepackage{amssymb}
\usepackage[margin=1.25in]{geometry}
\usepackage{microtype}
\usepackage{parskip}
\usepackage{graphicx}
\usepackage{booktabs}
\usepackage{float}
\usepackage{caption}
\usepackage{array}
\usepackage{xcolor}
\usepackage[numbers,sort&compress]{natbib}
\usepackage[colorlinks=true,linkcolor=blue,citecolor=blue,urlcolor=blue,
            pdftitle={A Perturbation Ladder for Sparse Autoencoder Dictionaries},
            pdfauthor={J. Melton}]{hyperref}
\usepackage[htt]{hyphenat}
\emergencystretch=2em

% ── Draft placeholder machinery ───────────────────────────────────────────────
\newcounter{pendingcount}
% Robust: \pending appears inside \caption, which is a moving argument -- a fragile
% macro there expands during .lof/.aux writing and crashes the engine.
\DeclareRobustCommand{\pending}[1]{\stepcounter{pendingcount}%
  \textcolor{red}{\textbf{[PENDING: #1]}}}
\makeatletter
\newcommand{\pendingtotal}{\textbf{?}}
\AtEndDocument{\immediate\write\@auxout{%
  \string\gdef\string\pendingtotal{\thependingcount}}}
\makeatother
\newcommand{\pendingnotice}{%
  \begin{center}\fcolorbox{red}{yellow!15}{\parbox{0.92\textwidth}{\centering
  \textbf{\textcolor{red}{DRAFT --- \pendingtotal\ unmeasured values remain.}}\\
  \small Every red bracket marks a number the experiment has not yet returned.
  No value in this document is estimated, extrapolated or illustrative:
  a number is either measured or visibly absent.}}\end{center}\medskip}

\title{\textbf{A Perturbation Ladder for Sparse Autoencoder\\Dictionaries: Seed,
Optimizer and Precision\\Against a Measured Ceiling}}

\author{
  J.\ Melton \\[4pt]
  American Code Labs \\
  \texttt{jmelton@americancode.org}
}

\date{September 2026 \\ \smallskip \small Preprint. Under review.}

\begin{document}
\maketitle
\pendingnotice

% ── Abstract ──────────────────────────────────────────────────────────────────
\begin{abstract}
This paper measures how much of a sparse autoencoder's learned dictionary survives
three perturbations---random seed, optimizer, and weight precision---using one matching
procedure and one denominator, so that the three are comparable to each other.

The seed rung supplies the reference the others are read against. Two TopK
autoencoders differing \emph{only} in random seed---identical model, layer, corpus,
architecture, optimizer and step count---agree on \textbf{7.9--8.3\%} of their
matchable features under Pearson fingerprints (8.9--10.0\% under cosine), recovered at
\textbf{108--112$\times$} a permutation null. Agreement is therefore strongly
significant and, simultaneously, far from total: retraining on identical data
reproduces under a tenth of a dictionary. Eight percent is the ceiling no perturbation
can exceed. Against it we measure optimizer (Adam versus Muon) and precision (bf16
versus locally converted 4-bit), and find \pending{headline: where precision and
optimizer land relative to the 8\% ceiling and the 0.18--0.33\% cross-architecture
floor}.

The precision rung addresses a question that arises because interpretability on
consumer hardware is usually done on quantized weights. \citet{duan2026quant} approach
it by holding a pretrained autoencoder fixed and applying it to full-precision and
quantized activations, reporting that roughly 62\% (Pythia-70M) and 51\% (Gemma-2-2B)
of active features survive INT6 at correlation above $0.9$, and identify comparison
against autoencoder \emph{training} variability as a complementary measurement their
design does not make. That is the measurement reported here. The two designs answer
different questions---whether a fixed dictionary's features still fire the same way,
versus whether the dictionary one would learn is the same dictionary---and their
percentages are quantities over different denominators that should not be compared
numerically.

Code, checkpoints and matching outputs are released; all runs are on a single Apple
M1 Max.
\end{abstract}

\tableofcontents
\newpage

% ─────────────────────────────────────────────────────────────────────────────
\section{Introduction}

Sparse autoencoders are trained on the residual stream of a model that a researcher
could afford to run. In practice that means quantized weights. Our own prior
cross-architecture study \citep{melton2026sae} had to disclose that one of its three
arms was quantized while the others were not, and could not say what that cost.

\subsection{Two Different Questions}
\label{sec:twoquestions}

There are two ways to ask whether quantization damages interpretability, and they are
not the same question.

\textbf{Frozen dictionary.} Take one trained SAE. Apply it to activations from the
full-precision model and from the quantized model. Ask whether each feature fires the
same way. This is \citet{duan2026quant}'s design, and it isolates the effect of weight
compression from anything to do with SAE training. It answers: \emph{if I already have
a dictionary, do its features still mean what they meant?}

\textbf{Relearned dictionary.} Train one SAE on full-precision activations and another,
identically, on quantized activations. Match the two dictionaries. This is our design.
It answers: \emph{if I do interpretability on a quantized model, do I discover the
same features?}

The second is the question a practitioner faces when they train an SAE on the only
model that fits their machine. It is also strictly harder to interpret, because
training an SAE twice does not give the same dictionary twice even with no perturbation
at all---which is exactly why it needs a baseline, and why the baseline is this paper's
central object.

\subsection{Why a Survival Number Needs a Ceiling}

\citeauthor{duan2026quant} report that 62.4\% of Pythia features survive INT6. Read
alone, the natural reaction is that roughly a third of interpretability is lost. But
neither their design nor ours has any meaning without knowing what perfect conditions
produce. They supply a null in the correct direction---an FP16-versus-FP16 check
returning correlation $1.0$ and zero damage---which establishes that their pipeline does
not manufacture drift, and under a frozen dictionary that check is the right one.

For the relearned-dictionary question no such trivial null exists, because rerunning
training is itself a perturbation. The ceiling has to be measured. We measure it.

\subsection{Contributions}

\begin{enumerate}
  \item \textbf{A measured ceiling for dictionary agreement.} Two autoencoders
    differing only in seed agree on 7.9--8.3\% of matchable features at
    108--112$\times$ the permutation null (\S\ref{sec:ceiling}). Any claim that a
    perturbation ``changes the features'' has to clear this.
  \item \textbf{A ladder rather than a point.} Seed, optimizer and precision measured
    on one instrument with one denominator, so the rungs are comparable to each other
    (\S\ref{sec:ladder}).
  \item \textbf{The complement \citeauthor{duan2026quant} identify as missing}: the
    training-variability comparison, on a locally converted deployed quantizer rather
    than simulated round-to-nearest (\S\ref{sec:method}).
  \item \textbf{An optimizer rung nobody has measured.} Whether the choice of optimizer
    changes the learned decomposition is, as far as we can determine, unexamined
    (\S\ref{sec:optimizer}).
  \item \textbf{Activation collection is a versioned artifact}
    (\S\ref{sec:notrepro}): released collection code is not sufficient to regenerate a
    released activation set, and arms collected under different collector revisions are
    not controlled.
\end{enumerate}

% ─────────────────────────────────────────────────────────────────────────────
\section{The Instrument}
\label{sec:instrument}

\subsection{Matching Procedure}

Features are compared by chunk-averaged activation fingerprints over a held-out corpus.
A feature pair is matched if each is the other's nearest neighbour---reciprocal
matching, which prevents a single popular feature from absorbing many partners---and if
similarity exceeds a threshold $\tau$.

$\tau$ is not chosen by hand. It is set at the 99th percentile of a permutation null in
which feature identity is shuffled, so that the expected number of false matches is
approximately 1\% of null matches. Enrichment over that null, rather than raw match
count, is the significance statistic.

Dead and low-support features are excluded before matching, and all percentages are
over \emph{matchable} features rather than dictionary size. This matters: a dictionary
of 16{,}384 with 3{,}300 dead and 900 low-support features has about 12{,}100 matchable
entries, and using the nominal size would deflate every figure by roughly a quarter.

\subsection{The Ceiling}
\label{sec:ceiling}

Three autoencoders were trained differing only in random seed---identical model
(Llama-3.2-3B, layer 14), corpus, architecture, dictionary size, $K$, optimizer, step
count and learning-rate schedule.

\begin{table}[H]
  \centering
  \caption{Seed control. Three pairwise comparisons among seeds 42, 123 and 456.
    ``\% matchable'' is calibrated matches over matchable features; enrichment is
    against the permutation null at $\tau = $ null $p_{99}$.}
  \label{tab:ceiling}
  \smallskip
  \begin{tabular}{llrrr}
    \toprule
    fingerprint & pair & $\tau$ & \% matchable & enrichment \\
    \midrule
    Pearson & 42 vs 123  & 0.9827 & 8.06 & 109.1$\times$ \\
    Pearson & 42 vs 456  & 0.9825 & 8.26 & 112.4$\times$ \\
    Pearson & 123 vs 456 & 0.9830 & 7.93 & 108.3$\times$ \\
    \midrule
    cosine  & 42 vs 123  & 0.9844 & 9.00 & 121.1$\times$ \\
    cosine  & 42 vs 456  & 0.9827 & 9.99 & 135.4$\times$ \\
    cosine  & 123 vs 456 & 0.9846 & 8.93 & 121.6$\times$ \\
    \bottomrule
  \end{tabular}
\end{table}

Two readings are both correct and both necessary. The agreement is \emph{highly
significant}---two orders of magnitude above chance, so the procedure detects shared
structure when shared structure exists. And the agreement is \emph{small}: more than
nine tenths of a dictionary is not reproduced by retraining on the same data with a
different seed.

\textbf{Consequences.} First, no perturbation can be shown to damage a dictionary by
more than reseeding already does; 8\% is the ceiling. Second, and less comfortably, a
published SAE dictionary is largely a sample from a distribution over dictionaries, and
claims about ``the features of layer $\ell$'' inherit that.

\subsection{The Floor}

The same procedure applied across architectures (Llama-3.2-3B, Pythia-1.4B,
Mistral-7B) returns 0.18--0.33\% of matchable features
\citep{melton2026sae}. Ceiling and floor together give every rung below a scale.

% ─────────────────────────────────────────────────────────────────────────────
\section{The Ladder}
\label{sec:ladder}

\begin{table}[H]
  \centering
  \caption{Perturbation rungs. Each arm differs from the bf16/Adam reference in exactly
    one respect; all share corpus, model, layer, dictionary size, $K$, batch, steps and
    seed 42.}
  \label{tab:rungs}
  \smallskip
  \begin{tabular}{llll}
    \toprule
    rung & what changes & agreement & vs ceiling \\
    \midrule
    seed      & random seed only        & 7.9--8.3\% & --- (defines it) \\
    optimizer & Adam $\to$ Muon         & \pending{} & \pending{} \\
    precision & bf16 $\to$ local 4-bit  & \pending{} & \pending{} \\
    \midrule
    floor     & architecture            & 0.18--0.33\% & \\
    \bottomrule
  \end{tabular}
\end{table}

\subsection{Precision}

The 4-bit arm is converted \emph{locally} from the same bf16 checkpoint with
\texttt{mlx\_lm.convert --q-bits 4}, not downloaded. A community upload cannot be
verified as to base revision or conversion settings, and a provenance gap in the one
variable under test would make the experiment unreadable. This also differs from
\citet{duan2026quant}, who simulate round-to-nearest quantization and list the absence
of a deployed quantizer as a limitation.

A sanity gate aborts the run if the bf16 and 4-bit activation files hash identically,
since that would mean the conversion had no effect and the comparison was vacuous.

\subsection{Optimizer}
\label{sec:optimizer}

Adam versus Muon \citep{jordan2024muon}, the orthogonalised-momentum optimizer whose
scaled results at Moonshot brought it to attention.

\textbf{The arms are matched on reconstruction quality, not on learning rate.} Adam and
Muon have natural learning rates two orders of magnitude apart, and Muon rescales
internally by $\sqrt{\text{rows}/\text{cols}}$. Holding the learning rate fixed would
pin a nuisance parameter and let quality float, so a low agreement score could mean
nothing more than that Muon was undertrained. A short probe selects the Muon learning
rate that matches Adam's FVE under an identical schedule, and the full run uses it.
This is also the control our own prior finding argues for: if standard quality metrics
cannot distinguish dictionaries that disagree about most of their content
\citep{melton2026sae}, then equal-quality dictionaries are precisely what should be
compared.

Two implementation requirements keep the arms separated by one variable. Muon's
Newton--Schulz step is defined for matrices, and MLX degrades to plain momentum SGD for
0D/1D parameters, so biases are routed to AdamW as that implementation's documentation
specifies. Muon's weight decay is set to zero rather than its default $0.01$, matching
the Adam arm, which applies none.

% ─────────────────────────────────────────────────────────────────────────────
\section{Method}
\label{sec:method}

\pending{full method table: model, layer, corpus, token budget, dictionary size, K,
steps, batch, schedule, hardware, wall-clock per arm}

\subsection{Activation Provenance}
\label{sec:notrepro}

Every arm is trained on activations collected as part of this work, from a corpus
regenerated and committed here, rather than on any previously collected set.

This is a requirement, not a preference. Activations previously collected for the
autoencoders of \citet{melton2026sae} cannot be regenerated from the released
collection script: the metadata schemas differ in field names, indicating the collector
was subsequently rewritten, and a hash over the leading $3.072 \times 10^{9}$ bytes
differs between the original and a fresh collection. An activation set that cannot be
regenerated cannot serve as a controlled arm, because the one variable under test would
share a file with an uncontrolled one.

The general form of the result is that \textbf{activation collection is a versioned
artifact, not a script}. Releasing the code that produced a set of activations is not
sufficient to reproduce them, and a pipeline that treats collection as reproducible by
construction will silently compare arms that were not collected the same way.

\section{Results}
\label{sec:results}

\begin{table}[H]
  \centering
  \caption{Agreement by rung. \pending{caption to state fingerprint metric and $\tau$}}
  \label{tab:results}
  \smallskip
  \begin{tabular}{lrrrr}
    \toprule
    comparison & $\tau$ & matches & \% matchable & enrichment \\
    \midrule
    bf16 vs 4-bit  & \pending{} & \pending{} & \pending{} & \pending{} \\
    Adam vs Muon   & \pending{} & \pending{} & \pending{} & \pending{} \\
    \bottomrule
  \end{tabular}
\end{table}

\subsection{Dictionary Quality Across Arms}

\pending{loss, FVE, L0, dead-feature count per arm. Prediction to be tested, not
assumed: quality metrics will be near-identical across arms regardless of how far the
dictionaries diverge, extending the metrics-blindness result of melton2026sae.}

% ─────────────────────────────────────────────────────────────────────────────
\section{Analysis}

\emph{To be written once Table~\ref{tab:results} is filled, and written once. The
measurement selects the argument; it is not a menu to hedge across.}

\textbf{If precision agreement approaches the 8\% ceiling}: quantization perturbs the
learned dictionary no more than reseeding does, and interpretability on 4-bit weights
is sound to the same degree that interpretability on any single SAE is sound---which is
a real qualification, but not one specific to quantization.

\textbf{If it falls well below}: features discovered on quantized models are
substantially not the features of the full-precision model, and a large amount of
consumer-hardware interpretability needs re-examination.

\textbf{If the optimizer rung moves as much as the precision rung}: the more
uncomfortable outcome, since optimizer choice is rarely reported as a threat to
validity at all.

% ─────────────────────────────────────────────────────────────────────────────
\section{Related Work}

\citet{duan2026quant} is the closest work and the reference point throughout;
\S\ref{sec:twoquestions} states the difference in question, and we do not present
competing numbers for the same quantity. \citet{bricken2023monosemanticity} and
\citet{cunningham2023sparse} establish the dictionary-learning setting;
\citet{gao2024topk} the TopK variant used here. \pending{additional coverage: SAE
reproducibility and feature-stability literature under vocabulary other than
``feature universality'', given that phrase-count surveys proved unreliable for this
purpose in companion work}.

% ─────────────────────────────────────────────────────────────────────────────
\section{Limitations}

\textbf{One model, one layer.} Llama-3.2-3B layer 14. Whether the ceiling is
8\% elsewhere is untested, and it is the number everything here is scaled by.

\textbf{One quantization scheme.} MLX 4-bit group quantization at one bit-width.
\citeauthor{duan2026quant} sweep INT4--INT8; we do not.

\textbf{One architecture family for the ladder.} The floor is cross-architecture but
the rungs are all Llama.

\textbf{Matching is one operationalisation of ``same feature''.} Reciprocal nearest
neighbour on activation fingerprints is not the only reasonable choice, and agreement
percentages are not invariant to it---as the Pearson/cosine spread in
Table~\ref{tab:ceiling} shows within this paper alone.

% ─────────────────────────────────────────────────────────────────────────────
\section{Reproducibility}

Code and data: \url{https://github.com/american-code/ResearchPapers}. Archived at
\pending{Zenodo concept DOI}.

\begin{center}
\begin{tabular}{ll}
\toprule
artifact & path \\
\midrule
run script          & \texttt{data/sae-runs/run\_quant\_interp.sh} \\
optimizer arm       & \texttt{data/sae-runs/run\_muon\_arm.sh} \\
trainer             & \texttt{data/sae-runs/train\_sae.py} \\
matching            & \texttt{data/sae-analysis/quant\_interp\_match.py} \\
seed control        & \texttt{data/seed-stability/} \\
\bottomrule
\end{tabular}
\end{center}

All runs on a single Apple M1 Max, 32\,GB unified memory, via MLX. Seed 42 throughout
except in the seed control.

\bibliographystyle{plainnat}
\bibliography{refs}

\end{document}
